
\section{Dam-break over a sand bed around a pier: the Zaragoza Kinect experiments}

Segovia-Burillo et al. \cite{SegoviaBurillo2026} released an 8~cm dam-break
from an $81 \times 157$~cm reservoir down a dry 6~m $\times$ 24~cm PVC flume
(Manning $n = 0.010$) whose bed, from 1.0 to 2.5~m past the gate, is a 5~cm
layer of sand (0.7 and 1.3~mm in equal parts, $\rho_s = 2650$~kg/m$^3$,
porosity 0.34) levelled with the PVC. A pier stands on the centreline 1.6~m
from the gate: a 3~cm cylinder (case P1) or a $3 \times 6.5$~cm rounded
rectangle (P2). The dam-break is repeated three times over the evolving
bed. A Kinect sensor over the pier reach ($x = 1.12$--2.07~m) records the
bed after each event and the water surface during each at 30~frames/s, on
a 1.4~mm grid; the data are open (Zenodo, DOI 10.5281/zenodo.17777387).
The flow over the sand is supercritical (Froude about 2), a few
centimetres deep, and lasts a few seconds; the transport is bedload. The
bed feature, the pier, is a few depths across and the scour hole many, so
the case sits inside the depth-averaged assumption in a way van Rijn's
trench does not, and it exercises the parts of the sediment model the
trench did not: the Exner coupling under a transient, a rigid base under a
finite layer, the angle of repose around a hole, and a wetting front over
an erodible bed.

\anuga{} meshes the reservoir and flume with the pier as a hole (reflective),
about 1~cm triangles over the sand and 5~mm around the pier, releases the
reservoir instantaneously (the gate opens in under 0.1~s, within the
Lauber--Hager criterion the authors check), and lets the flow out through
a transmissive boundary on the 4\,\% slope that follows the flat reach.
One bedload fraction of 1.0~mm (the geometric mean of the mix) moves under
Wong--Parker over a 5~cm erodible base with a 32$^\circ$ angle of repose;
the suspended exchange is off, as the paper treats the transport as
bedload. Manning $n$ on the sand is 0.015, the mean of the paper's values
for the two sizes. Each dam-break runs for 5~s of flow, after which the
water has left the pier reach; between events the reservoir is refilled and
the flume dried with the bed kept, as in the flume.

\begin{figure}[h]
\begin{center}
\includegraphics[width=0.95\textwidth]{setup.png}
\end{center}
\caption{The Zaragoza flume: an 8~cm dam-break runs over a 5~cm sand layer past a pier, in side view above and in plan below.}
\end{figure}

\IfFileExists{mesh.png}{
\begin{figure}[h]
\begin{center}
\includegraphics[width=0.95\textwidth]{mesh.png}
\end{center}
\caption{The mesh: about 1~cm triangles over the erodible sand and 5~mm
around the pier, which is cut out as a hole with reflective edges, on a
coarse reservoir and approach. The reservoir and the flume are one
continuous mesh, and the step at the gate line is a pair of wall segments.}
\end{figure}
}{}


Bedload is ramped off below a depth of two grain diameters, the
argument \texttt{min\_depth} of \texttt{set\_bedload}. Without it the film cells of the
wetting front, which carry the front's momentum and a large nominal shear,
dig steps into the sand as the front arrives and the time step collapses at
0.94~s, with or without the angle of repose and under DE0 or DE1.

The comparison is the bed change after each dam-break over the Kinect
window, block-averaged to 5~mm (the two \texttt{kinect\_*\_bed.csv} files
are that reduction of the Zenodo data, made by the script
\texttt{extract\_kinect.py}), plus the water surface along the centreline
during the first event. The Kinect bed change has an RMS of 3.6 (P1) and
4.0~mm (P2) about zero, most of it a ripple-scale texture the model does
not represent, so the provisional pass criterion is a bed-change RMS below
6~mm over the window after the first dam-break: \anuga{} gives 3.4 and
3.6~mm. What the case actually tests is the scour at the pier, and there
\anuga{} is short: over a ring 2--4~cm from the pier centre the flume
lowered the bed by 10.4~mm (P1) and 12.6~mm (P2) in the first event, the
model by 2.9 and 6.0; the deepest points are 32 and 53~mm against 16 and
12 (the 53 exceeds the sand thickness and is a Kinect artefact at the
pier's edge). The wake deposit is right in P1 (+2.3 against +3.0~mm) and
P2 (+2.2 against +2.4). Neither a 40$^\circ$ angle of repose, nor a mesh
twice as fine at the pier, nor turning the repose relaxation off moves the
ring scour by more than a millimetre, so the shortfall is in the transport
around the pier, not the discretisation: the horseshoe vortex that does
most of the digging at a pier is a three-dimensional structure a
depth-averaged model does not have, and the model's shear at the pier is
the uniform-flow one. The water surface in the window is not the pier's doing: with the pier
removed it is the same to 0.1~mm at 1.2 and 1.4~m, so the depth there is
set by the reservoir release and the flume, and the pier only adds a bow
wave (a 55~mm crest 20~cm upstream on a fixed bed at 4.7~s, absent once
the bed can scour). Against the Kinect the model's front is earlier and
blunter (21~mm at 1.2~m at 1.5~s against 9) and the later flow shallower
(37~mm at 1.2~m from 3.7~s against 44) and laterally flat, where the flume
shows cross-waves of $\pm 5$~mm from the contraction at the gate, a crest
on the centreline at 1.2~m and a trough at 1.4~m. Meshing the whole flume
at 1~cm does not produce them. The model's reservoir drops from 80 to
57~mm over the 5~s at a unit discharge of 0.027~m$^2$/s, between the
Ritter value for the depth (0.021) and the critical-contraction value
(0.038). Opening the gate over 0.15~s as a sill (\texttt{ZARAGOZA\_GATE\_TIME}),
refining the contraction at the gate (\texttt{ZARAGOZA\_A\_GATE}) and
raising the Manning $n$ of the sand were each tried: the gate delays the
front by 0.2~s and changes nothing after 2~s, the refinement changes
nothing, and $n = 0.025$ matches the later depth at 1.2~m to half a
millimetre while overshooting at 1.8~m, and none of the three moves the
ring scour past 4~mm. The shortfall is the transport at the pier itself.

\subsection*{The pier-scour correction}

The horseshoe vortex at a pier raises the bed shear to several times the
approach value within about a diameter, and a depth-averaged model has no
such structure. \anuga{} now carries a per-cell factor on the bed shear the
sediment sees (\texttt{set\_shear\_amplification}, [T-12]), and the case
sets it to $1 + A \exp[-(r - R)/(\ell R)]$ around the pier, $r$ the
distance from its axis and $R$ its half-width, chosen by
\texttt{ZARAGOZA\_PIER\_AMP} ($A$) and \texttt{ZARAGOZA\_PIER\_LEN} ($\ell$).
Its effect on P1 after the first dam-break:

\begin{center}
\footnotesize
\setlength{\tabcolsep}{3pt}
\begin{tabular}{p{3.0cm}rrrrr}
\hline
$A$, $\ell$ & RMS & ring 2--4~cm & deepest & wake & upstream \\
 & (mm) & (mm) & (mm) & (mm) & (mm) \\
\hline
measured & & $-10.4$ & 32 & +3.0 & $-1.4$ \\
0 (no correction) & 3.4 & $-2.9$ & 16 & +2.3 & +1.8 \\
2, 1 & 3.4 & $-3.6$ & & +3.6 & +0.7 \\
4, 1 & 3.4 & $-4.3$ & & +4.7 & $-0.2$ \\
8, 1 & 3.6 & $-5.5$ & & +6.3 & $-1.8$ \\
8, 1, repose off & 3.8 & $-4.2$ & 47 & +6.0 & $-2.1$ \\
8, 1, repose 45$^\circ$ & 3.6 & $-5.5$ & 31 & +6.4 & $-1.9$ \\
\textbf{4, 2} (the case) & 3.5 & $-5.0$ & 22 & +6.1 & $-2.0$ \\
8, 2 & 3.9 & $-7.0$ & 24 & +8.1 & $-4.8$ \\
4, 3 & 3.6 & $-5.1$ & 21 & +6.8 & $-2.5$ \\
8, 3 & 4.2 & $-7.4$ & 24 & +9.0 & $-5.4$ \\
8, 2, front and flanks only & 3.6 & $-1.5$ & 15 & +5.0 & $-2.4$ \\
\hline
\end{tabular}
\end{center}

The amplification digs, but not the way the flume did: the material it
lifts settles again within a pier length downstream, so the wake deposit
overshoots as the ring deepens, and the hole is narrow and deep against
the pier (a 45$^\circ$ repose gives the measured deepest point at $A = 8$)
where the flume's is broad and 10~mm deep over the whole ring. Confining
the amplification to the front and flanks, where the horseshoe vortex
lives, loses the flank scour altogether. The case runs with $A = 4$,
$\ell = 2$, which halves the ring deficit at no cost in RMS; a stronger
setting buys ring scour with wake and upstream error. What is missing is
not more shear but transport out of the wake: in the flume the wake
vortices keep the scoured sand moving, which is a suspension or a
turbulence effect the bedload closure cannot supply. P1 runs routinely and
P2 joins it under the long option.

\IfFileExists{bed_change_P1_event1.png}{
\begin{figure}
\begin{center}
\includegraphics[width=0.95\textwidth]{bed_change_P1_event1.png}
\end{center}
\caption{P1: bed change after the first dam-break, Kinect (top) and \anuga{} (bottom).}
\end{figure}
}{}

\IfFileExists{centreline_P1.png}{
\begin{figure}
\begin{center}
\includegraphics[width=0.95\textwidth]{centreline_P1.png}
\end{center}
\caption{P1: bed change along the flume centreline after each dam-break.}
\end{figure}
}{}

\IfFileExists{wse_P1.png}{
\begin{figure}
\begin{center}
\includegraphics[width=0.95\textwidth]{wse_P1.png}
\end{center}
\caption{P1: water surface above the initial bed along the centreline during the first dam-break.}
\end{figure}
}{}

\IfFileExists{bed_change_P2_event1.png}{
\begin{figure}
\begin{center}
\includegraphics[width=0.95\textwidth]{bed_change_P2_event1.png}
\end{center}
\caption{P2: bed change after the first dam-break, Kinect (top) and \anuga{} (bottom) (long run).}
\end{figure}
}{}

\endinput
